\scanpage{9}
% archive sheet number: 8
% lines 1
\[
 \operatorname{Ass}F=\handwritten{\bigcup_i}\operatorname{Ass}F_i
\]

% lines 2
(iii) $\operatorname{Ass}\handwritten{\coprod} F_i=\handwritten{\bigcup_i\operatorname{Ass}}F_i$ \qquad pour toute famille $(F_i)$ de modules qu c

% lines 3
Beweis klaar.

% lines 4–6
\uline{Proposition} 6.7. \quad Supposons $X$ loc noeth. Pour que la section $f\varepsilon I\underline O_X$ definisse une homothetie $F\quad F$ qui est \uline{injective}, il f et s que $X_f\quad\operatorname{Ass}F$, i.e. que $V(f)\quad\operatorname{Ass}F=\varnothing$.

% lines 7–8
En effet, on est ramene au cas affine (par exemple) o\handwritten{\`}{u} c'est dans Bourbaki. On en conclut:
% slanted
\margin{\handwritten{démonstrat\unsure{ion} directe}}

% lines 9–11
\uline{Corollaire} 6.8. \quad Soit $I$ un Ideal coherent sur $X$ loc noeth, $F$ un Module coherent sur $X$. Pour que $\underline I$ soit 'non diviseur de 0 dans $F$' il f et s que $V(I)\quad\operatorname{Ass}\unsure{\mathbb F|F}=\varnothing$.

% lines 12–13
On peut mettre ici si on veut IV 2.8.18. et 2.8.19. (bien qu'il soit possible qu'on ne les utilisera plus par la suite ...)

% lines 14–17
\uline{Proposition} 6.9. \quad Soient $X$ preschema loc noeth, $U$ un ouvert dense. Pour que $X$ soit reduit, il f et s que $U$ le soit, et que $X$ soit sans 'cycles premiers immerges'$^*$, ou encore que $X$ soit sans cycles premiers imm. et que l'on ait \quad ait long $\underline O_x=1$ pour $x$ point gen d'une comp irr de $X$.

% lines 18
(Cette notion * devrait etre definie, pour un Module, dans 6.1.).

% lines 19–26
Evidemment (i) \quad (ii) \quad (iii), d'autre part (iii) \quad (i) \deleted{puisque \unsure{d'}} d'apres le cas affine, bien connu \deleted{\unsure{(0} etant intersection des ideaux primaires \unsure{asso}} ou en notant qu'une sectionndu faisceau $\underline N$ des ideaux nilpotents a unnsupport qui ne peut contenir aucun point generique, donc s'il n'etait vide definirait par ses composantes irreductibles des cycles associes immerges. Il est peut-etre preferable d'essayer de donner \deleted{plutot} de telles demonstrations 'intuitives', plutot que de se reduire toujours a l'enonce connu dans le cas affine.

% lines 27–28
Nous allons moduler la notion de preschema reduit, en disant qu'un \uline{Module} qu coh $F$ \uline{est reduit} s'il est \uline{de type fini}, \uline{sans cycle premier associe}
